% ============================================================================
% Indonesian UI and structural strings for the One Math Book series.
% Loaded by styles/onemath.sty when \booklang is id.
%
% Standard Bahasa Indonesia, EYD spelling. Indonesian marks no plural on the
% noun, so every "plural" name below is deliberately identical to its singular
% — that is the correct form, not an oversight.
% ============================================================================

% Theorem / statement environment titles
\newcommand{\omnameDefinition}{Definisi}
\newcommand{\omnameTheorem}{Teorema}
\newcommand{\omnameProposition}{Proposisi}
\newcommand{\omnameLemma}{Lema}
\newcommand{\omnameCorollary}{Akibat}
\newcommand{\omnameMethod}{Metode}
\newcommand{\omnameExample}{Contoh}
\newcommand{\omnameNotation}{Notasi}
\newcommand{\omnameRemark}{Catatan}
\newcommand{\omnameExercise}{Latihan}
\newcommand{\omnameExercises}{Latihan}
\newcommand{\omnameProblem}{Soal}
\newcommand{\omnameProblems}{Soal}
\newcommand{\omnameProof}{Bukti}

% Admitted results and solutions appendix headers
\newcommand{\omadmittedtext}{Diterima tanpa bukti pada tingkat ini.}
% Used as "\omsolutionof~\cref{exo:...}": Indonesian needs no preposition here.
\newcommand{\omsolutionof}{Solusi}
\newcommand{\ompage}{halaman}
\newcommand{\omnameGotoSolution}{Solusi hlm.}

% Misc text macros
\newcommand{\st}{\text{ sedemikian sehingga }}

% Cover / title page
\newcommand{\bookmaintitle}{One Math Book}
\newcommand{\booktagline}{Satu Buku Matematika untuk Menguasai Semuanya}
\newcommand{\bookdescription}{Matematika dari taman kanak-kanak hingga akhir program sarjana}
\newcommand{\bookcontributors}{Para kontributor One Math Book}
\newcommand{\bookversionlabel}{Versi}

% Colophon
\newcommand{\omcolophonsource}{Kode sumber, laporan kesalahan dan kontribusi:}
\newcommand{\omcolophonlicense}{\textcopyright{} 2026 Benjamin Virrion.\\ Kecuali dinyatakan lain, buku ini dilisensikan di bawah Creative Commons CC BY-NC-SA 4.0:}

% Solutions appendix part title
\newcommand{\omsolutionspart}{Solusi Latihan}

% Part titles (Kelas 1–12; parallel to English Grade N)
\@namedef{omstr@part.grade1}{Kelas 1 (Usia 6--7 tahun)}
\@namedef{omstr@part.grade2}{Kelas 2 (Usia 7--8 tahun)}
\@namedef{omstr@part.grade3}{Kelas 3 (Usia 8--9 tahun)}
\@namedef{omstr@part.grade4}{Kelas 4 (Usia 9--10 tahun)}
\@namedef{omstr@part.grade5}{Kelas 5 (Usia 10--11 tahun)}
\@namedef{omstr@part.grade6}{Kelas 6 (Usia 11--12 tahun)}
\@namedef{omstr@part.grade7}{Kelas 7 (Usia 12--13 tahun)}
\@namedef{omstr@part.grade8}{Kelas 8 (Usia 13--14 tahun)}
\@namedef{omstr@part.grade9}{Kelas 9 (Usia 14--15 tahun)}
\@namedef{omstr@part.grade10}{Kelas 10 (Usia 15--16 tahun)}
\@namedef{omstr@part.grade11}{Kelas 11 (Usia 16--17 tahun)}
\@namedef{omstr@part.grade12}{Kelas 12 (Usia 17--18 tahun)}

% Part titles (Sarjana Tahun 1–3; parallel to English Bachelor Year N).
\@namedef{omstr@part.bachelor1}{Sarjana Tahun ke-1 (Usia 18--19 tahun)}
\@namedef{omstr@part.bachelor2}{Sarjana Tahun ke-2 (Usia 19--20 tahun)}
\@namedef{omstr@part.bachelor3}{Sarjana Tahun ke-3 (Usia 20--21 tahun)}

% Plural names + \omnameChapter, used by the \crefname declarations in
% onemath.sty. Identical to the singulars on purpose (see the header).
\newcommand{\omnameDefinitions}{Definisi}
\newcommand{\omnameTheorems}{Teorema}
\newcommand{\omnamePropositions}{Proposisi}
\newcommand{\omnameLemmas}{Lema}
\newcommand{\omnameCorollarys}{Akibat}
\newcommand{\omnameMethods}{Metode}
\newcommand{\omnameExamples}{Contoh}
\newcommand{\omnameNotations}{Notasi}
\newcommand{\omnameRemarks}{Catatan}
\newcommand{\omnameChapter}{Bab}
\newcommand{\omnameChapters}{Bab}

% LaTeX's own structural names. babel-indonesian sets the same words, but it
% is optional here (see onemath.sty), so declare them anyway.
\renewcommand{\chaptername}{Bab}
\renewcommand{\partname}{Bagian}
\renewcommand{\contentsname}{Daftar Isi}
\renewcommand{\appendixname}{Lampiran}
\renewcommand{\indexname}{Indeks}
